\chapter{Theory, Methodology, Materials \& Methods}

\section{Proposed System Architecture}
The system is organized into five stages: historical data backfill, live
daily ingestion, model training/retraining, a backend API layer, and a
two-page frontend application. Figure~\ref{fig:architecture} illustrates
the overall data and control flow.

\begin{figure}[h]
\centering
\begin{tikzpicture}[
    box/.style={draw, rounded corners, minimum width=3.4cm, minimum height=1cm, align=center, font=\small},
    arr/.style={-Latex, thick}
]
\node[box] (kaggle) at (0,6) {Kaggle Historical\\Backfill (one-time)};
\node[box] (live) at (5,6) {Agmarknet Live API\\(daily)};
\node[box] (db) at (2.5,4.5) {State-partitioned\\Database};
\node[box] (train) at (2.5,3) {Retraining: one\\Prophet model per\\(commodity, mandi)};
\node[box] (registry) at (2.5,1.5) {Model Registry};
\node[box] (api) at (2.5,0) {FastAPI Backend\\(+ state aggregation)};
\node[box] (frontend) at (2.5,-1.5) {React Frontend\\(Home + Prediction pages)};

\draw[arr] (kaggle) -- (db);
\draw[arr] (live) -- (db);
\draw[arr] (db) -- (train);
\draw[arr] (train) -- (registry);
\draw[arr] (registry) -- (api);
\draw[arr] (api) -- (frontend);
\end{tikzpicture}
\caption{Proposed system architecture: data sources, storage, model
training/registry, API, and frontend.}
\label{fig:architecture}
\end{figure}

An important design decision, discussed further in Section~3.4, is that
\textbf{storage and the user interface are organized by state}, while the
\textbf{forecasting model itself remains at (commodity, mandi)
granularity} -- state-level views are computed as an aggregation over
underlying mandi-level forecasts rather than being separately modelled.

\section{Data Pipeline}
\subsection{Historical Backfill}
Prophet requires several months of history to reliably detect weekly and
yearly seasonality; this cannot be obtained from live daily collection
alone within a single academic year. The system therefore performs a
one-time historical backfill from the Kaggle "Daily Market Prices of
Commodity India (2001--2026)" dataset -- specifically its 2024 and 2025
slices -- before daily live collection begins. This dataset is compiled
by the same original maintainer as the live Agmarknet API dataset used
for daily ingestion, giving one consistent data lineage across historical
and live collection.

\subsection{State and Crop Selection}
The three states and six crops per state (Table~\ref{tab:cropselection})
were finalized by ranking commodities within the 2024+2025 backfill data
by total reported record count per commodity per state (aggregated
across all mandis and days), with mandi count and day count examined
alongside as supporting context. For Tamil Nadu, whose top commodities
were closely clustered in volume, rankings were computed on combined
2024+2025 totals rather than a single year, since single-year rankings
shifted noticeably between the two years; Uttar Pradesh and Maharashtra
produced the same top-6 ranking independently in both years and did not
require this adjustment.

\begin{longtable}{@{}p{2.8cm}p{10.5cm}@{}}
\toprule
\textbf{State} & \textbf{Top 6 crops (ranked)} \\
\midrule
\endhead
Tamil Nadu & Coconut; Bhindi (Ladies Finger); Green Chilli; Bottle gourd;
Snakeguard; Onion (edged out Banana-Green by a narrow margin, ranked on
combined 2024+2025 data) \\
Uttar Pradesh & Potato; Onion; Tomato; Wheat; Brinjal; Green Chilli
(stable ranking independently in both 2024 and 2025) \\
Maharashtra & Wheat; Bengal Gram; Soyabean; Arhar (Tur/Red Gram); Onion;
Jowar (Sorghum) (stable ranking independently in both 2024 and 2025) \\
\bottomrule
\caption{Finalized state and top-6-crop selection, ranked by total
reported record count on the 2024+2025 backfill data.}
\label{tab:cropselection}
\end{longtable}

\subsection{Live Ingestion}
A scheduled job (GitHub Actions cron) fetches the previous day's records
from the Agmarknet API (resource ID
\texttt{9ef84268-d588-465a-a308-a864a43d0070}) and appends them to the
same state-partitioned tables used by the historical backfill, so that
the model's training set grows continuously.

\subsection{Data Cleaning Considerations}
Preliminary analysis of a sample Agmarknet export surfaced two practical
issues that the ingestion pipeline must handle:
\begin{itemize}
    \item \textbf{Field-name encoding.} Exported column names contain an
    XML-encoding artefact (e.g.\ \texttt{Min\_x0020\_Price}) rather than
    the clean \texttt{min\_price} field name returned by the JSON API;
    the pipeline must normalize this during ingestion.
    \item \textbf{Variety/grade duplication.} A single (market, commodity,
    date) combination can have multiple rows if more than one
    variety/grade of that commodity was traded that day (confirmed
    empirically, e.g. multiple \textit{Arecanut} rows at the same market
    on the same day with different modal prices). The pipeline filters to
    a consistent variety/grade per (commodity, mandi) series before
    constructing the Prophet training set.
\end{itemize}

\section{Machine Learning Design}
\subsection{Task Formulation}
The core forecasting task is \textbf{time-series regression}, not
classification: the model predicts a continuous value (the next-period
modal price), using \texttt{modal\_price} as the target variable (the
price at which the largest transacted quantity occurred that day, per
Agmarknet convention) rather than the min/max range. The farmer-facing
Sell/Hold/Stable label is a rule-based decision layer applied on top of
the continuous forecast, not a directly trained classifier -- this
preserves the underlying magnitude/confidence information that a bare
label would discard, and mirrors how real trading/advisory systems are
typically structured.

\subsection{Model Choice}
\textbf{Facebook Prophet} (an additive decomposition model: trend +
weekly seasonality + yearly seasonality) is used as the primary
forecasting model, trained independently per (commodity, mandi) pair. A
naive/moving-average baseline is trained alongside it for comparison,
evaluated using Mean Absolute Error (MAE), to provide an evidence-based
justification for the modelling choice rather than an assumed one.

\subsection{State-Level Aggregation (not a separate model)}
The Prediction page's state-level summary view ("select a state, see its
top 6 crops with percentage change") is computed by aggregating the
outputs of the existing per-mandi models -- e.g.\ averaging the forecast
percentage change across all mandis reporting a given commodity within
the selected state -- rather than training a new, coarser model. This
preserves per-mandi precision for the underlying comparison feature while
still providing a fast, simple state-level overview.

\section{Technology Stack}
Table~\ref{tab:techstack} summarizes the finalized technology choices and
their justification.

\begin{longtable}{@{}p{2.8cm}p{3.8cm}p{6.5cm}@{}}
\toprule
\textbf{Layer} & \textbf{Technology} & \textbf{Justification} \\
\midrule
\endhead
Historical backfill & Kaggle "Daily Market Prices of Commodity India" (2024+2025 parquet) & Same lineage/maintainer as live API dataset; bulk-filterable across states/crops \\
Live data source & data.gov.in Agmarknet API & Free, structured, government-maintained \\
Forecasting model & Facebook Prophet & Handles small per-mandi series well; minimal tuning \\
Baseline model & Naive / moving average (scikit-learn) & Enables an evidence-based accuracy comparison \\
Backend & FastAPI (Python) & Lightweight, async, auto-documented REST API \\
Database & SQLite (dev) / Supabase Postgres (prod) & Free hosting; state-partitioned schema \\
Frontend & React (Vite) & Required for interactive, clickable India map \\
Mapping & react-simple-maps + TopoJSON & Standard approach for clickable choropleth/state maps \\
Charting & recharts & React-native price/forecast visualization \\
Scheduling & GitHub Actions & Free CI/CD-based cron scheduling \\
Deployment & Vercel/Netlify (frontend), Render (backend), Supabase (DB) & All free-tier, standard for student deployment \\
\bottomrule
\caption{Finalized technology stack.}
\label{tab:techstack}
\end{longtable}

\section{Feasibility Discussion}
\subsection{Data Feasibility}
Both the live (Agmarknet API) and historical (Kaggle-compiled Agmarknet)
data sources have been confirmed to be freely accessible without
institutional gatekeeping,
unlike domains such as healthcare (patient-data access restrictions) or
finance (proprietary transaction data), which was a key factor in
selecting this domain over alternatives originally considered.

\subsection{Technical Feasibility}
Each component of the proposed stack (Prophet, FastAPI, React) is
well-documented and widely used in comparable applications, keeping the
implementation risk manageable for a team building most of this tooling
for the first time in a practical setting.

\subsection{Timeline Feasibility}
The project scope (2--3 states, six commodities per state, roughly
15--20 mandis per state) was deliberately bounded to fit a single
academic-year capstone timeline, with a clear path to widening coverage
later without architectural change (Section~3.1).

\subsection{Ethical and Regulatory Considerations}
As the system uses only publicly published government price data and
does not process any personal, medical, or financial data, it does not
raise the consent, privacy, or regulatory concerns that a healthcare- or
finance-domain project would.
